\documentclass[11pt]{article}
\usepackage[utf8]{inputenc}
\usepackage[T1]{fontenc}
\usepackage{lmodern}
\usepackage[margin=1in]{geometry}
\usepackage{csquotes}
\usepackage[style=authoryear,backend=biber,sorting=nyt,maxbibnames=99,giveninits=true,uniquename=init]{biblatex}
\usepackage[hidelinks]{hyperref}
\addbibresource{references.bib}
\begin{document}
\section*{Controlled Drainage and Irrigation Water Management}

\subsection{Overview}
Integrating controlled drainage (CD) with irrigation water management is repeatedly framed less as a settled practice than as a research frontier: the recurring call is to move away from passive, one-way ``unilateral'' subsurface drainage toward systems that give farmers active, location-specific control over the water table and over the coupled irrigation--drainage water balance \parencite{Ayars2006b,Kesicka2023b,Ritzema2023}. Across arid, semiarid, humid, and irrigated command-area settings, authors converge on a small set of open problems: (i) developing design criteria, structures, and decision-support tools for \emph{active} rather than passive management \parencite{Ayars2006b,Youssef2023,Saadat2017}; (ii) automating and data-driving water-balance control through sensors and models \parencite{Youssef2023,Ikonen2023,Isomäki2024}; (iii) coordinating irrigation with drainage adjustments, including subirrigation embedded in regional water systems \parencite{Wahba2005,deWit2022,deWit2024}; (iv) sustaining drainage-water reuse and salinity control while protecting downstream water quality \parencite{WESTCOT1988,Sharma1998,Ayars2015}; (v) adapting design to changing crop water needs, climate extremes, and irrigation modernization \parencite{Rui2024,Ritzema2007,Ziaeifar2025}; and (vi) resolving the institutional, economic, and participatory barriers that condition adoption \parencite{Ritzema2011,Castellano2019,Qureshi2008}. A cross-cutting message is that many recommendations remain qualitative and site-specific, and that long-term, multi-scale evidence is still lacking. Kesicka \textcite{Kesicka2023b} frames replacing unilateral drainage with methods of controlling drainage outflows as a current global challenge, and continued research is broadly expected to keep refining drainage techniques to optimize agricultural water use.

\subsection{From passive to active, adaptive management}
A foundational research direction is the transition from ``no or passive management to active management with defined water management objectives'' \parencite{Ayars2006b}. \textcite{Ayars2006b} argue that structures for controlling water-table position should be \emph{required} in new drainage design and considered for retrofitting existing systems where practical, and that there is a pressing need to develop new design criteria and management methods for controlled drainage to support sustainable irrigated agriculture---a claim they note requires long-term field studies, possibly combined with modeling, for validation. Similar reasoning appears in system-design recommendations to include structures for controlling groundwater-table (GWT) position in new subsurface drainage designs \parencite{Dukhovny2018}, to incorporate ``more options for control'' so farmers can optimize integrated water management \parencite{Ritzema2023}, and to expand ``changeable controlled drainage networks'' across arable lands \parencite{Mahmoud2021}. \textcite{Strock2011} anticipate that producers may eventually view controlled and reduced drainage as soil- and water-management \emph{tools} rather than mitigation practices, and argue that fine-tuning drainage intensity should be a priority in designing all new systems.

The quantitative payoff of active management remains an open question. \textcite{Watkins2024} call for research on the crop-yield benefits achievable by actively managing the subsurface CD outlet during extremely wet periods---hypothesizing that active management could minimize wet-stress on roots and raise yields---and separately for study of how shallow drainage combined with CD affects drainage and nutrient loading and downstream water quality. \textcite{Saadat2017} highlight the need to examine the cost--benefit trade-off of active management, and \textcite{Helmers2012} suggest that yield reductions observed under CD could likely be mitigated with greater management of the system, while cautioning that installation cost, producer management burden, and the acreage of applicability must be weighed before broad implementation. Seasonally differentiated control is another explicit hypothesis: \textcite{Asgari2021} propose free drainage for quick waterlogging control in normal and wet seasons, CD in dry seasons to alleviate deficit stress, and possibly supplemental surface ditches and supplemental irrigation to manage the extremes.

\subsection{Sensors, automation, and decision-support systems}
Several groups identify automation and data-driven control as the next technological step. \textcite{Youssef2023} call for a new generation of drainage control structures and decision-support tools for active management of CD during the growing season, for sensor-based, data-driven water-management systems on poorly drained soils with frequently shallow water tables, and more broadly for technologically advanced systems that adapt to changing weather patterns. \textcite{Ayars2015} likewise identify decision-support systems for adaptive control of integrated irrigation and drainage systems as a need. In Finland, \textcite{Ikonen2023} lay out a concrete research agenda: improved dynamic modeling and data assimilation of the field water balance, a feasible water-balance control-problem formulation (cost function and constraints), fusing control automation with the physical drainage control-well infrastructure, and gaining practical experience from end-users and experimental peatland agriculture. \textcite{Isomäki2024} argue for data-driven water-management systems and appropriate control schedules for combined controlled subsurface drainage and open-ditch damming, noting that proactive management of control structures according to meteorological conditions can increase the probability of keeping water-table depth in the optimal range. Automatic water-depth control is also proposed for aerobic rice, where \textcite{Arif2021} introduce a conceptual sheet-pipe automatic water-depth control system that still needs field testing, alongside an opportunity to examine the ET$_c$--production relationship for aerobic rice. Earlier work anticipated these directions, describing farm decision aids and management modeling systems---slowly being developed---that would incorporate water quantity and quality, crop growth, management alternatives, and economic evaluation \parencite{THOMAS1995}.

\subsection{Integrating irrigation with drainage and subirrigation}
Coordinating irrigation adjustments with drainage control is a persistent open problem. \textcite{Wahba2005} stress that when CD is implemented, appropriate reductions in irrigation application must follow, requiring coordination and training among irrigation authorities, drainage authorities, and farmers. A recurring recommendation is to improve integration between irrigation and drainage networks and to use CD to help farmers optimize on-farm water management \parencite{Ritzema2011,Ritzema2016}. \textcite{Dukhovny2018} argue that controlled drainage (CTD) requires a genuinely integrated irrigation-and-drainage water-management approach, with careful site selection (favoring homogeneous soil texture with little sand or sand--gravel near the control structure) and careful monitoring of all water-balance components.

Subirrigation embedded in regional water systems is a distinct and largely unresolved frontier. \textcite{deWit2022} argue that regional-scale water needs for subirrigation should be known before responsible incorporation into regional water management, framing correct implementation of CD with subirrigation as an integral, catchment-wide question; they call for spatial and temporal insight into water availability at each location, optimized pumping strategies to minimize water use, and cautious use of alternative resources such as industrial or treated wastewater only where quality is sufficient and side effects are limited. The companion modeling study \parencite{deWit2024} shows how the SWAP model could identify suitable field sites and forecast the hydrological consequences of subirrigation before implementation, explore effects across crops and crop characteristics, and examine drainage design, water-supply design, soil type, resistance to downward seepage, and surface-water management under varying meteorology; it also suggests optimizing water supply by starting later, supplying water only when crop transpiration is potentially reduced, and coordinating surface-water levels to avoid drainage losses to ditches. Longer-term studies are needed to disentangle whether soil-property changes under drainage-subirrigation arise from the practice itself or from interactions with irrigation water quality, and whether drip or sprinkler methods interact with soil properties \parencite{Kaur2023}. Related integrative directions include continuing to combine water-table control (via CD and subirrigation) with agrochemical applications to minimize nutrient and pesticide losses \parencite{BENGTSON1995}, and pairing shallow drainage with wider perennial-crop strips along drainageways \parencite{Schilling2007}.

\subsection{Design criteria for changing crops, climate, and irrigation modernization}
A broad theme is that design criteria must be revisited for evolving conditions. \textcite{Rui2024} argue that future drainage design should account for changing crop water needs, evolving weather, and recent engineering advances, suggesting more efficient---ideally controllable---systems to improve long-term productivity of poorly drained soils. Design criteria should be revisited to meet environmental concerns \parencite{Akram2013}, and drainage design and water management should be optimized in the face of climate change through adaptive strategies such as adjusting planting schedules, improving water productivity, and enhancing drainage to mitigate waterlogging \parencite{Ziaeifar2025}. \textcite{Ritzema2023} recommend explicitly incorporating the impact of irrigation modernization into drainage design to account for changing soil-water and salinity regimes. Ageing systems, new crop varieties, crop diversification, water scarcity, deteriorating water quality, and wet/dry climate extremes are all cited as drivers of continued drainage research \parencite{Ritzema2007}, and changes in drainage-water quality are expected to require changes in drain-system design and installation \parencite{Ritzema2006}. Further research is needed to serve emerging and least-developed countries and to cope with climate, land-use change, and drainage-water quantity/quality requirements \parencite{Ritzema2011}.

Design intensity itself is an open research question. \textcite{Dialameh2025} note that where adjusting drain spacing is infeasible, other water-management practices may be needed to mitigate increased drought severity, and that future design decisions may need to account for management and economic factors beyond those modeled. \textcite{Singh2009} caution that reduced drain spacing would require increased capacity across the entire drainage network---a factor for future network design---and would heighten the emphasis on best management practices to reduce nutrient loss. \textcite{Castellano2019} advocate a systems approach integrating drainage, crop--soil processes, and nutrient-loss-reduction practices to design drained cropping systems that mitigate and adapt to climate change, and call for identifying long-term cost savings (e.g., greater fertilizer nitrogen-use efficiency) that could offset the higher installation costs of altered drain depth and spacing. \textcite{Strock2011} recommend detailed analysis with DRAINMOD or another drainage model to develop optimum drainage rates and drain spacing across U.S. locations and soils. Fundamental design questions persist about which soil parameters are most critical and what confidence intervals are required for each \parencite{Ayars2015}, and about proper design guidelines for specific geologic conditions such as the Mid-Atlantic states \parencite{SHIRMOHAMMADI1995}. The optimum design is repeatedly framed as a cost--benefit balance \parencite{HARDJOAMIDJOJO1982}, and irrigation is emphasized as necessary to reduce deficit soil-water conditions where drainage alone cannot raise yields \parencite{HARDJOAMIDJOJO1982}.

\subsection{Drainage-water reuse, salinity, and water-quality frontiers}
Managing salt and reusing drainage water remain heavily research-dependent. \textcite{Jia2007} frame understanding salt-and-water-balance dynamics as the first step toward better management of irrigated agriculture, providing evidence to target and fairly evaluate field water-management practices and improve regional water-use efficiency. \textcite{Ayars2015} enumerate needs to develop management approaches for subsurface drainage in arid areas, to reduce crop salt-tolerance and groundwater-management concepts to a manageable producer system, to quantify drainage-water-treatment efficiency (loading, flow rates, lifespan, maintenance, design), to refine salt-management strategies for integrated systems, and to develop and evaluate alternatives for managing, remediating, and removing nitrogen from drainage water.

Reuse-focused research spans agronomy, scheduling, and safety. \textcite{WESTCOT1988} list needs for site-specific load--flow relationships for trace elements, expanded knowledge of germination and early-seedling salinity tolerance, ratings of crops' use of saline groundwater, improved irrigation scheduling that lets farmers meet crop ET from a portion of the high water table while minimizing upward salt flow, wider commercial applicability of direct drainage-water reuse, better understanding of bioaccumulation and trace-element immobilization in evaporation basins, and caution in reusing salts until safety is better known. Field guidance likewise remains provisional: substituting drainage water up to 9~dS/m for wheat at post-germination stages where non-saline canal water is scarce, then reverting to non-saline canal water for one or two years after 6--7 years of drainage-water use, with careful case-by-case evaluation \parencite{Sharma1998}. Keeping drainage effluent as low as possible by improving irrigation efficiency and blending fresh and drainage water, plus disposal options such as farm ponds for storing and evaporating drainage water, are proposed \parencite{Datta2002}, as is minimizing effluent by improving irrigation efficiency and farm-level reuse while maintaining soil health \parencite{Qureshi2008}. Autumn or winter irrigation combined with subsurface pipe drainage is suggested for leaching excess soil salt \parencite{Dou2022}, seasonal groundwater undulation is deemed necessary to prevent salt accumulation \parencite{Mirlas2021}, and proper irrigation schemes---shifting from traditional to drip irrigation---are recommended to sustain growth while flushing salts \parencite{Mirlas2021}. Root-zone salinity and the water table should be kept low, at least in early growth stages \parencite{Singh2019}. Non-conventional treatment via engineered in-stream and off-stream wetlands is being explored at pilot scales \parencite{AbdelDayem2007}, and drainage water outflow is argued to require treatment before discharge to surface or groundwater \parencite{Filipovic2015}. Emerging soil-remediation directions include exploring acidic biochar in saline-alkaline soils and its synergy with subsurface drainage, plus plant-remediation desalination and other collaborative measures \parencite{Liu2025f}. Reclamation innovations include harvesting subsurface-drainage water in dryland/semiarid sites for livestock, windbreaks, gardens, or salt leaching, and establishing alfalfa in saline-seep areas fitted with engineered drainage wells \parencite{Steppuhn2017}. Progressive-flooding and progressive-leaching methodologies still require experimental validation and extension to finer-textured soils \parencite{Youngs2000}. Salinity control is argued to belong within a basin-wide salt-balance context \parencite{Qureshi2008}, and comprehensive assessment of waterlogging and salinity together (rather than separately) is recommended to yield more reliable findings \parencite{Valipour2014}.

\subsection{Nutrient, pesticide, and water-quality research needs}
Water-quality outcomes of controlled and water-table-management systems remain incompletely quantified. \textcite{FAUSEY1995} identify a need to evaluate the full potential of CD and water-table management for water quality, spanning short- and long-term effects on the fate and transport of agricultural chemicals, soil properties, productivity and trafficability, surface/subsurface water and sediment movement, root-zone biological dynamics, and plant development across crops; they also flag CD for animal-waste disposal on cropland, winter denitrification of residual soil nitrogen via CD, interfacing drainage with constructed or enhanced wetlands, and---as a high priority---a comprehensive water-table-management model incorporating solute transport, with results folded into practical management guides. \textcite{RITTER1995} add needs to quantify preferential flow in drainage-water pollution, develop field- and watershed-scale pollution-control modeling, determine year-round nitrate and pesticide leaching, gain more field experience with subirrigation and CD, document tillage effects on water quality, interface agrichemicals with livestock manures where tile drainage is used, and investigate cover crops for water-quality management. Continued pesticide degradation/movement studies and improved simulation of chemical losses under high water tables are being pursued \parencite{THOMAS1995}, and monitoring of drainage impacts on nutrient and pesticide transport---treating drainage as part of whole-farm or watershed management---is recommended \parencite{SHIRMOHAMMADI1995}. At watershed scale, controlling or slowing surface runoff and tile flow, with drainage-water-management devices and denitrification bioreactors at tile outlets, is proposed to cut nitrate loads \parencite{Li2015}, and strategies addressing both hydrology and nitrogen supply, focused at the drainage-district scale, are argued necessary for meaningful downstream benefits \parencite{Ikenberry2014}. Optimizing irrigation is crucial to minimize nitrogen leaching (e.g., from turfgrass) \parencite{Filipovic2015}, and coupling CD with conservation practices that reduce surface runoff is recommended \parencite{BinShokrana2023}. Reduced tillage should be combined with efficient subsurface drainage to minimize surface runoff \parencite{Koskiaho2002}, and CD is offered as one way to meet both crop-production and water-quality goals \parencite{Kladivko2005}. Nutrient replenishment under the drainage--irrigation--cropping system---including zinc fertilization---is needed for optimum yields \parencite{Gul2013}, and understanding nutrient-loss mechanisms during waterlogging and the benefits of nutrient application is an open question \parencite{Manik2019}.

\subsection{Modeling, monitoring, and long-term evaluation}
Authors repeatedly call for stronger modeling, longer records, and better monitoring. Long-term studies are broadly needed to validate active-management approaches \parencite{Ayars2006b}, and the long-term effects of CD---especially in arid and semiarid regions---remain understudied \parencite{Wang2020c}. \textcite{Wang2020c} further urge defining soils by internationally accepted standards (or fuller descriptions), investigating CD from a climate standpoint, combining future designs with proposed underlying mechanisms for more reliable evaluation, and classifying situations to understand CD's contribution to total nitrogen in drainage water. Watershed-scale water-table-management research is needed to develop best management practices and select the best water-management scenario, and more flexible CD schemes may alter impacts on tile drainage, surface runoff, and shallow-water-table days \parencite{Golmohammadi2016}. \textcite{Kim2023} present the SWAT-tile model as opening opportunities for water-quality research and tile-depth installation, and as a viable strategy for managing flow and nitrate transport at field and watershed scales. Steady/transient simulations of steady recharge are framed as an initial step toward broader analyses that could propose conceptual modifications to traditional design for water-quality control \parencite{GRISMER1993}. Regional modeling calls include extending simulation periods to capture dynamic equilibrium and long-term sustainability, completing calibration of groundwater flow-and-transport models, and refining field data on irrigation efficiencies, canal seepage, drainage, and return flows \parencite{Triana2010}. Environmental monitoring and decision-support systems should be designed at project onset to work hand-in-hand with drainage operations, built bottom-up and documented progressively \parencite{Quinn2009}. Additional model-supported directions include using a different field-scale hydrologic model to evaluate CD's yield effect \parencite{BinShokrana2023} and exploring interactive tools such as DRAINFRAME \parencite{Gopalakrishnan2007}. Remote sensing to detect drainage and salinity needs---via drones and satellite imagery \parencite{Yannopoulos2020}, and refined, upscaled remote-sensing assessment of waterlogged/saline areas \parencite{Gopalakrishnan2007}---is another data frontier. Reviewing prior research and projects is urged to avoid trial-and-error and reduce failure risk in future rehabilitation/construction \parencite{Valipour2014}.

\subsection{Crop, soil, and agronomic research directions}
Refining the biophysical targets of control is an active frontier. The required water-table depth must be refined by crop, soil type, and season, and it must be verified that shallow drains can control root-zone salinity over the long term before farmer implementation \parencite{Christen2001b}. Root-development responses to tile drainage---and consequent effects on plant water uptake during drought---need further study \parencite{Nash2015b}. Further work is needed on drainage and desalination in unsaturated soil and on the design/engineering of tile drainage under drip irrigation \parencite{Chen2021}. Winter crops that leave soils bare during hot summers warrant more investigation, and longer research is needed on CD in sugarcane \parencite{Akram2013}. On soils with high perched water tables, progress depends on understanding factors limiting crop growth and developing improved strategies, potentially including mole drainage or increasing drainable porosity by loosening soil above the drain \parencite{Muirhead1996}; strategies to reduce flow to shallow subsurface drains and more sustainable practices are also needed \parencite{Muirhead1996}. Soil-management options such as subsoiling, deep-rooted cropping, and drainage trenches are viable short-term solutions \parencite{Guedessou2021}. \textcite{Manik2019} outline a waterlogging-tolerance agenda: increasing soil-profile de-watering through soil-improvement and agronomic strategies, greater specificity of management--environment and management--management interactions, discovering new non-oxygen-associated QTLs, testing their effectiveness alone and combined under multiple constraints, evaluating agronomic/yield/quality effects, and developing management packages for varieties with diverse tolerance genes---plus cost/benefit comparison of different drainage strategies. Salt-tolerant germplasm, groundwater-abstraction regulation, and crop diversification (e.g., salt-tolerant medicinal and aromatic species) are recommended \parencite{Qureshi2008}, as are salt-tolerant crops such as eucalyptus or phreatophytes and drainage installation where leaching with poor-quality water is infeasible \parencite{Sarwar2001}. Supplementary measures---gypsum application, salt-resistant varieties, and irrigation-efficiency improvement---are needed to enhance subsurface-drainage benefits \parencite{Ritzema2008b}. Pressurized irrigation and reduced irrigation amounts are recommended to lessen waterlogging and drainage needs \parencite{Khorramian2022}, more research is needed on leaching requirements for modern irrigation in heavy clay soil \parencite{Mahmoud2021}, and alternatives such as surface irrigation and improved fertility management alongside proper drainage should be weighed for claypan soils \parencite{MOSTAGHIMI1989}.

\subsection{Economics, institutions, and adoption}
A large share of the frontier is socio-economic and institutional. Economic analysis is repeatedly flagged as future work \parencite{Kladivko2005,Saadat2017,Manik2019}, and installation cost, producer management burden, and applicable acreage must be assessed before broad implementation \parencite{Helmers2012}. A broader economic analysis is urged to reveal opportunities for improving water management, reducing drainage costs, and enhancing land/water quality, including the opportunity cost of drainage-investment funds, the scope to shrink drainage facilities through better water management, and pricing/allocation policies that motivate improvement \parencite{Wichelns2002}. The overall economic component of water-quality-management systems needs to be determined and incorporated into design \parencite{THOMAS1995}, and identifying long-term cost savings to offset higher installation costs, plus incentives and cost-share programs that couple drainage upgrades with edge-of-field nitrate-loss reduction, is recommended \parencite{Castellano2019}. Cost reduction is itself a research target: local manufacture of PVC pipes and synthetic envelopes and further labor-cost reductions \parencite{Datta2002}; and new low-cost drainage techniques, cheaper installation technologies, new drainage materials (pipes, envelopes), and national standards \parencite{Yannopoulos2020}.

Institutional and participatory needs are pervasive. Recommendations include shifting from standardized to flexible, location-specific solutions reflecting stakeholder preferences \parencite{Ritzema2011,Ritzema2016,Ritzema2023}; balancing top-down against bottom-up through participatory management across planning, design, installation, operation, and maintenance \parencite{Ritzema2016}; improving farmer participation throughout the drainage-infrastructure lifecycle \parencite{Ritzema2023}; and better identifying stakeholders' needs, preferences, and willingness to contribute while linking their tacit knowledge with researchers' explicit knowledge \parencite{Ritzema2011}. Capacity development---integrating research, education, and advisory services---is emphasized \parencite{Ritzema2016}, along with the substantial education and policy change needed to lower adoption barriers \parencite{Castellano2019}, public/institutional education programs on drainage and water-table management \parencite{SHIRMOHAMMADI1995}, and addressing negative perceptions that conflate agricultural drainage with wetland drainage before the application base can expand \parencite{THOMAS1995}. Scientists and policy institutions should help farmers adopt innovations through financial support and knowledge exchange \parencite{Kesicka2023b}. Governance recommendations include developing procedural/organizational frameworks and national or regional support groups, working groups for diverse agro-hydro-climatic regions, and multi-departmental coordination committees in command areas \parencite{Singh2019c}; enhanced inter-agency coordination for water allocation \parencite{Qureshi2008}; a dedicated state-level Drainage Division and national drainage institution, standardization of components, and promotion of private-sector reclamation and local SSD installation \parencite{Gopalakrishnan2007}; drainage policy papers with time-bound reclamation action plans \parencite{Ritzema2008b,Valipour2014}; and appropriate institutional arrangements for farmers' participation \parencite{Ritzema2008b}. Adaptive regulatory frameworks that promote local irrigation/drainage flexibility and grower innovation, communicate successes, recognize failures early, and disseminate information are recommended \parencite{Quinn2009}. Programs to create awareness and promote subsidies and micro-finance \parencite{Satyanarayana2007}, to convince farmers to clear drainage blockages and maintain systems \parencite{Manjunatha2004}, and to ensure participatory operation and maintenance for sustainability \parencite{Gopalakrishnan2007} round out the adoption agenda. Continued technical and financial support is deemed essential to sustain irrigated agriculture \parencite{Qureshi2008}, research in irrigated agriculture is ``not yet over'' \parencite{Ritzema2007}, and a participatory approach is demanded by changing socio-economic conditions \parencite{Ritzema2007}. Practical framing choices---outsourcing/privatizing drainage activities \parencite{AbdelDayem2007}, budgeting for certified local consultants to support decentralized design \parencite{Ritzema2023}, and switching from gravity to pumped drainage to enable reuse and shrink open-drain dimensions \parencite{Ritzema2023}---are also proposed. Large-scale drainage investment tied to adjusted irrigation planning is viewed as unavoidable in the long run \parencite{Sarwar2001}, and accumulated reclamation experience is expected to inform future developments worldwide \parencite{TeKrony2004}.

\subsection{Implementation, maintenance, and scaling}
Finally, several frontier items concern getting systems built, maintained, and scaled. CD could be propagated on a large scale in irrigation command areas \parencite{Karegoudar2019}, and adopting CD, drain-water reuse, and standardized components is called for \parencite{Gopalakrishnan2007}. For large-scale implementation, trencher and trenchless drainage machines may be considered \parencite{Ritzema2008b}. Later-stage management and maintenance are highlighted as important, with sludge removal from drainage pipes identified as the chief issue in later utilization \parencite{Yan2013b}; upgrading and carefully maintaining existing tile drainage---and exercising more care in construction and maintenance---is recommended for better efficiency \parencite{AbdElMageed2009}. Efficiency-oriented site actions include better canal maintenance, gate operations, and application techniques, converting gravity delivery to closed (piped) delivery, widely implementing drip and sprinkler networks, preventing erosion/sedimentation in storage structures, ensuring good subsurface-drainage management against salinity and waterlogging, and avoiding deep-groundwater use to prevent saltwater intrusion \parencite{Kanber2019}. Modeling procedures can pre-screen appropriate field sites and forecast hydrological consequences before implementing subirrigation \parencite{deWit2024}. Sustainable effluent disposal must ultimately be resolved \parencite{Datta2002}, low-cost and easily understood CD options should be pursued within an integrated water-saving approach \parencite{Wahba2005}, and profitable, feasible water-management alternatives for problem soils should be selected on a site-specific basis \parencite{MOSTAGHIMI1989,Sharma1998}.

\printbibliography
\end{document}
